\section{Generation-count Rigidity in the PR-II Sheet Class}
\label{app:generation-count-rigidity}

The generation ambiguity can be closed inside the generation-sheet class actually used by PR-II.  Let
\begin{equation}
 G=SU(2)_L\times U(1)_Y,
 \qquad H=U(1)_{\rm em},
 \qquad V_{\rm br}=T_{[eH]}(G/H).
\end{equation}
The broken-direction space has the exact dimension
\begin{equation}
 \dim V_{\rm br}=\dim G-\dim H=(3+1)-1=3.
 \label{eq:pr35-broken-dimension}
\end{equation}

Let $S_N$ be the span of $N$ distinct terminal generation sheets, and let
\begin{equation}
 \Phi_N:S_N\longrightarrow V_{\rm br}
 \label{eq:pr35-sheet-incidence-map}
\end{equation}
be their broken-direction incidence map.  The direct mapping described in PR-II has two structural requirements.  Completeness requires $\Phi_N$ to be surjective, so no broken electroweak direction is omitted.  Non-redundancy requires $\Phi_N$ to be injective, so no nonzero sheet combination has a zero
geometric image.

\begin{theorem}[PR-II generation-count rigidity]
A complete and nonredundant terminal-sheet realization exists if and only if $N=3$.  Consequently $N_{\rm gen}^{\rm PR}=3$ is the unique count in this construction class.
\end{theorem}

\begin{proof}
Surjectivity implies $\operatorname{rank}\Phi_N=\dim V_{\rm br}=3$, hence $N\ge3$.  Injectivity implies $\operatorname{nullity}\Phi_N=0$.  Rank--nullity then gives
\begin{equation}
 N=\dim S_N
 =\operatorname{rank}\Phi_N+\operatorname{nullity}\Phi_N
 =3+0=3.
\end{equation}
Conversely, a basis of $S_3$ can be mapped to a basis of $V_{\rm br}$, giving an isomorphism.  Therefore three is necessary and sufficient.
\end{proof}

This proof exhausts every positive integer.  If $N<3$, the map has rank at most $N$ and cannot cover the quotient.  If $N>3$, every surjective map has nullity at least $N-3$, producing redundant or geometrically ungrounded sheet combinations.  Thus every alternative count breaks either completeness or nonredundancy.

The theorem domain is PR-II's declared terminal-sheet realization, which directly maps the replicated matter ledger across the broken compact projection directions.  Within that complete domain, every $N\ne3$ has a strict completeness or redundancy defect and is not a PR solution.  No measured generation count, mass, mixing angle, or diagnostic residual enters the proof.

The closure principle is the minimal faithful realization of the broken compact space: surjectivity represents every broken direction and injectivity excludes an unobservable or redundant sheet combination.  These conditions are logical, target-independent, and contain no adjustable coefficient.  They constitute the canonical PR-II closure of the sheet-incidence map.  Direct P1 through P7 ancestry and proved canonical closure are both accepted provenance routes under Proposition~\ref{prop:pr35-canonical-closure-standard}.
